\section{A model of review rounds}
\label{sec:model}

\subsection{Documents and edits}

We view a manuscript as a finite sequence of source characters over an
alphabet $\Sigma$, $\doc \in \Sigma^{*}$, of length $n = \abs{\doc}$. An
\emph{edit} replaces a half-open range of the source by new text.

\begin{definition}[Edit]
An edit is a triple $e = (a, b, w)$ with $0 \le a \le b \le n$ and
$w \in \Sigma^{*}$. Applying $e$ to $\doc$ gives
\begin{equation}
  \doc \apply e = \doc_{[0,a)} \, w \, \doc_{[b,n)},
  \label{eq:apply}
\end{equation}
whose length is $n + \delta(e)$ with \emph{shift} $\delta(e) = \abs{w} - (b - a)$.
\end{definition}

A task proposes a \emph{revision} $\edit = \{e_1, \dots, e_m\}$, a finite set
of edits against the same base document. Papilot groups nearby edits into
\emph{changes}, the units the author decides on.

\begin{definition}[Compatibility]
Two edits $e = (a, b, w)$ and $e' = (a', b', w')$ are \emph{compatible},
written $e \perp e'$, when their ranges are disjoint and not adjacent:
\begin{equation}
  e \perp e' \iff b < a' \ \text{or}\ b' < a .
  \label{eq:compatible}
\end{equation}
A revision is compatible when its edits are pairwise compatible.
\end{definition}

For a compatible revision we apply edits from right to left, so that earlier
offsets remain valid:
\begin{equation}
  \doc \apply \edit
  = \bigl(\cdots\bigl((\doc \apply e_{\sigma(1)}) \apply e_{\sigma(2)}\bigr)
    \cdots\bigr) \apply e_{\sigma(m)},
  \qquad a_{\sigma(1)} > a_{\sigma(2)} > \dots > a_{\sigma(m)} .
  \label{eq:apply-revision}
\end{equation}

\begin{proposition}[Compatible edits commute]
\label{prop:commute}
If $e \perp e'$, then applying them right to left in either listing gives the
same document, and the result of \cref{eq:apply-revision} does not depend on
how the remaining edits of a compatible revision are ordered.
\end{proposition}

\begin{proof}
Assume without loss of generality $b < a'$. Applying $e'$ first changes only
positions at or after $a'$, so the range $[a, b)$ and every offset before it are
unchanged, and applying $e$ afterwards gives
$\doc_{[0,a)}\, w\, \doc_{[b,a')}\, w'\, \doc_{[b',n)}$. Applying $e$ first
shifts every later offset by $\delta(e)$; rewriting $e'$ as
$(a'+\delta(e), b'+\delta(e), w')$ yields the same string. Induction on the
number of edits, ordering them by start offset, gives the general statement.
\end{proof}

\subsection{Rebasing edits}

When one task's edit $e$ has been accepted, another task's edit $e'$, written
against the old document, must be moved onto the new one. The
\emph{transformation} of $e'$ past $e$ is
\begin{equation}
  \tau(e', e) =
  \begin{cases}
    (a', b', w'), & b' < a,\\
    \bigl(a' + \delta(e),\ b' + \delta(e),\ w'\bigr), & b < a',\\
    \bot, & \text{otherwise,}
  \end{cases}
  \label{eq:transform}
\end{equation}
where $\bot$ signals a conflict that the author must resolve.

\begin{lemma}[Convergence of rebased edits]
\label{lem:tp1}
If $e \perp e'$, then
\begin{equation}
  \doc \apply e \apply \tau(e', e) = \doc \apply e' \apply \tau(e, e') .
  \label{eq:tp1}
\end{equation}
\end{lemma}

\begin{proof}
If $b < a'$, then $\tau(e, e') = e$ and $\tau(e', e)$ is $e'$ shifted by
$\delta(e)$, so the left side is the second computation in the proof of
\cref{prop:commute} and the right side the first. The case $b' < a$ is
symmetric.
\end{proof}

\Cref{eq:tp1} is the convergence property known from operational
transformation~\cite{ellis1989}. Papilot needs only the compatible case: an
incompatible pair is never merged automatically.

\subsection{Decisions and rounds}

Let $U$ be the set of changes in a round. A \emph{decision state} is a map
\begin{equation}
  \theta : U \to \{\textsf{pending}, \textsf{accepted}, \textsf{rejected},
  \textsf{commented}\},
\end{equation}
and the accepted part of a revision is
$\edit^{\theta} = \{ e \in \edit : \theta(u(e)) = \textsf{accepted} \}$.
The manuscript after round $t$ is
\begin{equation}
  \doc_{t+1} =
  \begin{cases}
    \doc_t \apply \edit_t^{\theta_t}, & \text{if } \edit_t^{\theta_t}
      \text{ applies cleanly to } \doc_t,\\[2pt]
    \doc_t, & \text{otherwise, and the task is asked for a new proposal.}
  \end{cases}
  \label{eq:round}
\end{equation}
Rejected and commented changes return to the agent as feedback bound to their
round, file and lines; the agent must not reintroduce a rejected edit unless it
is explicitly asked to.

\paragraph{Actions and undo.}
An \emph{action} $\alpha = (\iota, u, d)$ carries an identifier $\iota$, a
change $u$ and a decision $d$. Applying it to a state gives
$\theta \cdot \alpha = \theta[u \mapsto d]$, and the history is a stack
$\alpha_1, \dots, \alpha_k$. Two properties make decisions safe over an
unreliable network:
\begin{align}
  (\theta \cdot \alpha) \cdot \alpha &= \theta \cdot \alpha
    && \text{(idempotence: a retried write changes nothing)},
    \label{eq:idempotent}\\
  \operatorname{undo}(\theta_0 \cdot \alpha_1 \cdots \alpha_k)
    &= \theta_0 \cdot \alpha_1 \cdots \alpha_{k-1}
    && \text{(undo reverses exactly the last action)}.
    \label{eq:undo}
\end{align}
An undo request names the action it reverses, so a retried undo cannot reverse
a different one.

\subsection{How many rounds?}

Suppose that each change is accepted in a round independently with probability
$p \in (0,1]$, and otherwise revised and proposed again. The number of rounds
$R$ until a single change is accepted is geometric,
\begin{equation}
  \Prob(R = r) = (1-p)^{r-1} p, \qquad
  \E[R] = \frac{1}{p}, \qquad
  \Var(R) = \frac{1-p}{p^{2}} .
\end{equation}
For $m$ independent changes the task ends when the last one is accepted, after
$R_{\max} = \max_i R_i$ rounds.

\begin{proposition}[Expected rounds]
\label{prop:rounds}
\begin{equation}
  \E[R_{\max}]
  = \sum_{k=1}^{m} (-1)^{k+1} \binom{m}{k} \frac{1}{1 - (1-p)^{k}} .
  \label{eq:rounds}
\end{equation}
\end{proposition}

The proof is in \cref{app:proofs}. For small $p$ the sum behaves like
$\frac{1}{p} H_m$, where $H_m = \sum_{k=1}^{m} 1/k \approx \ln m + \gamma$ is
the $m$-th harmonic number: doubling the number of changes adds only about
$(\ln 2)/p$ rounds.

\begin{corollary}[Rounds with confidence]
\label{cor:rounds}
For every $\eta \in (0,1)$, the task ends within
\begin{equation}
  r_{\eta} = \ceil*{ \frac{\ln(m/\eta)}{-\ln(1-p)} }
  \label{eq:rounds-confidence}
\end{equation}
rounds with probability at least $1 - \eta$.
\end{corollary}

\begin{proof}
By the union bound, $\Prob(R_{\max} > r) \le m\, \Prob(R > r) = m (1-p)^{r}$,
which is at most $\eta$ as soon as $r \ge \ln(m/\eta) / (-\ln(1-p))$.
\end{proof}

\begin{example}
With $p = 0.8$ and $m = 5$, \cref{eq:rounds} gives
$\E[R_{\max}] \approx 1.91$ and \cref{eq:rounds-confidence} with $\eta = 0.05$
gives $r_{\eta} = 3$: most tasks finish in one or two rounds, and almost all
within three.
\end{example}
